\chapter*{Preface}
\addcontentsline{toc}{chapter}{Preface}

Book 10 closed the geometric tier of this series with a functor, $F_t$, representing a single lamphron-lamphrodyne oscillation, and the categorical periodicity $F_{t+\tau} = F_t \circ F_\tau$ that functor satisfied. The present book takes that closing construction as its starting point rather than as an isolated curiosity: everything the geometric tier built, connections, metrics, curvature tensors, entropy functionals, is here reorganized as a category, with curvature itself becoming syntax, the structure that allows reality to compose itself. Geometry needed grammar, and this book supplies it.

\chapter*{Diagrammatic Overview}
\addcontentsline{toc}{chapter}{Diagrammatic Overview}

The geometric hierarchy built across Books 6 through 10 maps onto categorical structure directly: manifolds become objects, curvature-preserving maps between them become morphisms, the entropy and curvature constructions of those earlier books become functors out of the resulting category, and the naturality conditions this book develops express exactly the thermodynamic consistency those books maintained by more piecemeal argument.

\chapter*{Reading Note}
\addcontentsline{toc}{chapter}{Reading Note}

Composition replaces expansion as this series' central organizing operation, and smoothness becomes coherence: every claim in the geometric tier about a manifold smoothing under entropy flow has a categorical restatement in this book as a claim about some diagram commuting, and the reader should treat these restatements as literal translations rather than as loose analogies.

\part{The Objects and Morphisms of RSVP}

\chapter{Objects: Thermodynamic Manifolds}

The category $\mathbf{RSVPField}$ has as its objects tuples $X = (M, g, \PhiField, \vField, \SEntropy)$, a manifold equipped with the metric Book 6 constructed and the field triple that metric was built from, subject to the fixed plenum measure $\mathrm{vol}(M) = \text{const}$ already required by the non-expanding plenum constraint this series has maintained since that book. Objects of this category are read here as sites of persistence, locations at which a particular configuration of curvature, entropy, and flow is realized and sustained rather than merely passing through.

Examples developed in this chapter, a cognitive surface in the sense of Book 9, an ecological basin, and a social attractor, each instantiate this object structure in a different domain, illustrating that $\mathbf{RSVPField}$'s objects are not restricted to any single scale of application: whatever satisfies the defining tuple, at whatever scale, is a legitimate object of this category.

